\begin{statement}{E.45}{Proved}{The complete $A_6$ kernel and outer-coordinate tests retain every old fibre}
Let $B=C_2*C_3$, $F=\langle x,y\rangle$, $N$ and $K_6=N'N^6$
be the discrete groups with their literal generators from E.1 and E.24.
The action of the Ferrari Nielsen generators on the epimorphisms
$F_2\twoheadrightarrow A_6$ has $53$ orbits under
$\operatorname{Aut}(A_6)$, with Nielsen orbit sizes
$12,15,10,16$.  Retain either the permutation of these kernels,
or that permutation together with every transporter in
$\operatorname{Out}(A_6)\simeq C_2^2$.
These give specified quotient maps
\[
 B_4\longrightarrow T_1\longrightarrow T_0,
 \qquad B\xrightarrow{\rho_\epsilon}G_\epsilon
 \quad(\epsilon=0,1),
 \tag{E.45.1}
\]
where $\rho_\epsilon$ retains all six Ferrari evaluations and
all six cosets of $F$ in $B$.
Put $C_\epsilon=\rho_\epsilon(N')$ and
$S_m=B/(F'F^m)$.  Then
\[
 \begin{gathered}
 C_\epsilon=\rho_\epsilon(F')
            =\rho_\epsilon(K_6\cap F'),
 \qquad C_\epsilon'=C_\epsilon,\\
 G_\epsilon/C_\epsilon=S_{30},\qquad
 G_\epsilon/\rho_\epsilon(K_6)=S_6,
 \qquad |S_{30}|=5400,\quad |S_6|=216.
 \end{gathered}
 \tag{E.45.2}
\]
The two equalities with $S_m$ identify the actual maps from $B$.
Both projections $T_1\to T_0$ and $G_1\to G_0$ have
elementary abelian kernel of order $2^{35}$.

Every one of the $864$ retained points of E.24 extends to these
common source and four-strand targets.  The assertion holds for
every profinite multiplier congruent to $1$ modulo $72$, with
the appropriate retained E.24 point.  For the source the actual
common image and the corrected image are
\[
                 D_{23}\times_{S_6}G_\epsilon;
 \tag{E.45.3}
\]
for the four-strand target they are
\[
                 T_{23}\times_{C_2}T_\epsilon,
 \tag{E.45.4}
\]
where the two maps to $C_2$ are braid exponent modulo two.
All these maps are surjective and satisfy the complete pentagon,
braid and full-twist equations.

For multiplier one the correction can simultaneously preserve
any full boundary fibres of E.42.  Thus all $432$ points of E.24
coexist with these two $A_6$ tests, the complete specified
two- and three-adic class-two boundary fibres, and the identity
on every full relative pro-$p$ boundary quotient for $p\ge5$.
For cofinal old common sources retaining $P_{20}$ from E.43,
the actual common source and its corrected image are
\[
                 G_{\rm old}\times_{S_{30}}G_\epsilon.
 \tag{E.45.5}
\]
This assertion concerns the displayed family of quotients.
\end{statement}
\begin{proof}
We give the finite model and each comparison used in the lifting argument.
We use the complete $A_6$ and $\operatorname{Aut}(A_6)$ tables of D.31,
and independently check every one of the $1440$ automorphisms on
both generator edges of the entire $360$-vertex Cayley graph.
The action partitions all $129600$ ordered pairs into $124$ orbits.
Testing generation on one representative of every orbit gives
$76320$ generating pairs.  The automorphism action on generating
pairs is free, so there are exactly $76320/1440=53$ epimorphism
kernels.  All pairs, orbit representatives and transporters are retained.

For representatives $(g_j,h_j)$, the two literal Nielsen operations are
\[
 (g,h)\longmapsto(g,hg),\qquad
 (g,h)\longmapsto(gh^{-1},h).
 \tag{E.45.6}
\]
For each operation the program records the unique permutation $p$
of the $53$ kernels and actual automorphisms $c_j$ carrying
$(g_{p(j)},h_{p(j)})$ to the displayed pair.  It checks the entire
automorphism-valued braid relation before taking any quotient.
The multiplication is
\[
 (p,c)(q,d)=\bigl(q\circ p,(c_jd_{p(j)})_j\bigr).
 \tag{E.45.7}
\]
Keeping only $p$ gives $T_0$; keeping $p$ and the two outer bits
of every $c_j$ gives $T_1$.  The latter acts faithfully on
the $212$ pairs $(j,e)$ with $e\in C_2^2$, by
$(j,e)\mapsto(p(j),e+\overline c_j)$.
The $53$-point action is its specified quotient.
The two actual Ferrari generators $\xi_1,\xi_2$ map to identity
in these targets, and consequently their three Artin generators
are $(a,b,a)$.  Direct multiplication verifies all three
$B_4$ braid relations.  This defines the first maps in
\textup{(E.45.1)}.

This construction retains the universal epikernel coordinates
of \cite[Lemma 2.1]{DE_Guillot2015}: the free group evaluates
onto $A_6^{53}$, and the Nielsen substitutions act on this
product by the displayed permutation and actual automorphisms.
The maps just defined compose the Ferrari action on that product
with its kernel-permutation or outer-coordinate projection.
No identification of the finite automorphism group computed in
\cite[Theorems 4.9 and 5.2]{DE_Guillot2015} with our source image
is used.

Let $\theta_i:F\to T_\epsilon$, $1\le i\le6$, be the
six pure-word evaluations printed in E.17, now using these
literal $a,b,a$.  Let $\tau_q$ be the retained Schreier section
of $B/F\simeq S_3$.  The source homomorphism is the induced map
whose $36$ coordinates at $b\in B$ are
\[
 \theta_i\bigl(\tau_q b\tau_{q\overline b}^{-1}\bigr),
       \qquad q\in B/F,\quad 1\le i\le6,
 \tag{E.45.8}
\]
together with the permutation $q\mapsto q\overline b$.
The intervening word lies in $F$; the exact free-graph
Schreier map converts it to a word in $x,y$.
Multiplication in the induced group follows by cancellation
of the adjacent section words.  Thus its kernel is normal
in $B$, and it retains all six evaluations and their conjugates.
It extends continuously to the profinite completion.

The resulting faithful permutation degrees are $1908$ and $7632$.
We reduce them to $432$ and $798$, respectively, by deleting
only repeated transitive $B$-sets.  For every deleted orbit the
certificate gives a bijection to a retained orbit and checks
that it intertwines both original generators $s,r$ pointwise.
These maps prove equality of the two kernels, so the compressed
actions are faithful representations of the same $G_\epsilon$.

Write
\[
 c_0=2^{40}3^{21}5^{10}7^6 11^3 13^2,
       \qquad c_1=2^{35}c_0.
 \tag{E.45.9}
\]
Complete permutation-group calculations give
\[
\begin{gathered}
\begin{array}{c|rrrr}
 &G_\epsilon&\rho_\epsilon(F)&\rho_\epsilon(F')&\rho_\epsilon(N)\\ \hline
 |\ \cdot\ |/c_\epsilon&5400&900&1&75
\end{array}\\[4pt]
\begin{array}{c|rrr}
 &\rho_\epsilon(N')&\rho_\epsilon(K_6)&\rho_\epsilon(K_6\cap F')\\ \hline
 |\ \cdot\ |/c_\epsilon&1&25&1
\end{array}
\end{gathered}
 \tag{E.45.10}
\]
Here is the completeness of the subgroup calculation.
All $97$ literal Schreier generators of $N$ are retained,
and the program verifies that the four selected generators
generate their entire finite image.  It computes the derived
subgroup, including the full normal closure of the commutators,
then adjoins every actual sixth power for $K_6$.
For $K_6\cap F'$ it uses the $95$ integral relations among
the two ordinary exponent sums of the $97$ generators.
The integer relation matrix has rank $95$ and Smith diagonal
$(1^{95})$, and its product with the exponent-sum matrix is zero.
It is therefore the full saturated integral kernel.
Modulo $\rho_\epsilon(N')$ the products of the sixth powers
commute, so these $95$ products give precisely the image of
$K_6\cap F'$.  Each is already in $\rho_\epsilon(N')$.
The derived subgroup of this last image has the same order,
proving its perfectness.  The full derived subgroup of the
two-generator $F$ image has the same order and contains it,
proving the first line of \textup{(E.45.2)}.

The literal images of $x,y$ in both sources have order $30$.
Since they generate $F$ and its image modulo $C_\epsilon$
is abelian of order $900$, the specified map
\[
 (\mathbb Z/30)^2\longrightarrow\rho_\epsilon(F)/C_\epsilon,
                 \qquad (u,v)\longmapsto x^uy^v C_\epsilon
 \tag{E.45.11}
\]
is an isomorphism.  Consequently $B\to G_\epsilon/C_\epsilon$
factors through $S_{30}$; both groups have order $5400$, so
the factorization is the claimed isomorphism.  The image of
$K_6$ contains $x^6,y^6$, as checked directly, and has order
$25c_\epsilon$.  The corresponding quotient is therefore
the actual $S_6$ quotient.  These comparisons use ordinary
free exponent sums, with no change of their normalization.

The upper and lower models are related by forgetting the outer
bits in every coordinate of \textup{(E.45.8)}.  The orders in
\textup{(E.45.9)}--\textup{(E.45.10)} show that its source kernel
has order $2^{35}$.  In that kernel every permutation coordinate
is identity, so multiplication of the remaining bits is addition
in an elementary abelian two-group.  The same calculation on
the four-strand targets gives
\[
       |T_\epsilon|=2c_\epsilon,
       \qquad |T_\epsilon'|=|T_\epsilon''|=c_\epsilon.
 \tag{E.45.12}
\]
Thus $T_\epsilon'$ is perfect, and its maximal solvable quotient
is the actual braid-exponent map $T_\epsilon\to C_2$.
The target projection also has elementary abelian kernel of
order $2^{35}$.  The full twist $(aba)^4$ is identity in both
targets; $(ab)^3$ has order one in $T_0$ and two in $T_1$.

We now prove the full-fibre assertion.
Fix a retained E.24 word $f_0\in\widehat F'$ and its multiplier
$\lambda\equiv1\pmod{72}$.  Choose an actual Galois element
with cyclotomic character $\lambda$, and write its standard
Ihara word as $f_\sigma\in\widehat F'$.
The cyclotomic character is onto $\widehat{\mathbb Z}^{\times}$,
so this choice is available for every stated multiplier.
By \textup{(E.45.2)},
\[
       \rho_\epsilon(f_0)^{-1}\rho_\epsilon(f_\sigma)
                                 \in C_\epsilon.
 \tag{E.45.13}
\]
The joint image of $N'$ in the old finite source $D_{23}$
and $C_\epsilon$ projects onto each factor.  The first image
is solvable and the second is perfect.  Every common quotient
is both solvable and perfect and is therefore trivial.
Goursat's lemma gives the direct product, providing a word
$k\in N'$ with identity old image and the value in
\textup{(E.45.13)} in the new source.  Then $f=f_0k$
retains the complete old point and agrees with $f_\sigma$
in all $36$ new source coordinates.  In particular all six
four-strand evaluations agree.  The actual Galois action,
in the braid convention of E.38, proves all the asserted
equations and surjectivity on the new factors.  Every finite
power uses the same chosen profinite $\lambda$; this argument
requires no restriction on its new five-primary residue.
For $\lambda=1$ we may choose $\sigma=1$, so all new evaluations
of the corrected $f$ are literally identity.

The exact generation comparisons are as follows.
Every common quotient of $D_{23}$ and $G_\epsilon$ kills
$C_\epsilon$, since $D_{23}$ is solvable.  It also kills
$\rho_\epsilon(K_6)/C_\epsilon\simeq C_5^2$, since
$D_{23}$ has only two- and three-primary factors.
It is therefore a quotient of $S_6$.  Conversely $S_6$ is
a specified quotient of both groups: the old source retains
$G_6$, and $K_6\subset F'F^6$.
Thus $S_6$ is their maximal common quotient with these maps
from $B$.  Goursat proves \textup{(E.45.3)} for both the original
source and the corrected image.  On $S_6$ the prescribed
maps agree, since the multiplier is $1$ modulo six and
the word lies in $F'$.
Similarly, \textup{(E.45.12)} shows that every common quotient
of $T_{23}$ and $T_\epsilon$ factors through $C_2$.
On the binary Ferrari factor the action on the three nonzero
vectors of $\mathbb F_2^2$ sends each Artin generator to a
transposition, giving this actual common quotient.
Odd multipliers and conjugation preserve it.  This proves
\textup{(E.45.4)} and the asserted joint four-strand generation.

For the simultaneous boundary assertion let $H$ and the closed
derived terms $\widehat H^{[r]}$ be as in E.42--E.43, and put
$R_H=\bigcap_r\widehat H^{[r]}$.
The exact $H$ image in $S_{30}$ is
$3(\mathbb Z/30)^2\simeq C_{10}^2$ under
\textup{(E.45.11)}.  In particular
\[
 B\longrightarrow P_{20}
       =B/\bigl(\gamma_3(H)H^{20}\bigr)
                         \longrightarrow S_{30}
 \tag{E.45.14}
\]
is the literal quotient factorization, and $S_{30}$ is the
maximal solvable quotient of $G_\epsilon$.
Perfectness gives $\rho_\epsilon(\widehat N^{[r]})=C_\epsilon$
for every $r\ge1$.  The inclusions
$\widehat N^{[r+1]}\subset\widehat H^{[r]}$ and compactness of
the nested fibres over each element of $C_\epsilon$ prove
\[
                  \rho_\epsilon|_{R_H}:R_H\twoheadrightarrow C_\epsilon.
 \tag{E.45.15}
\]
The image is contained in $C_\epsilon$ because $S_{30}$ is
solvable.  Choose this preimage to cancel the new value of the
E.42 word $f_0$.  It changes none of the finite solvable
$H$-images, hence preserves the entire stated boundary family
simultaneously.  All its new evaluations are identity.

Every finite old common source in E.42 is solvable: its
$H$ image is a finite product of finite $p$-groups and its
quotient $B/H$ has order $54$.
For finite old common sources containing $P_{20}$, the maximal
common quotient is now exactly $S_{30}$, by
\textup{(E.45.14)} and maximal solvability.
Both component maps generate, and they agree on this specified
quotient.  Goursat gives \textup{(E.45.5)}.  Such old sources are
cofinal in the family of E.42, since adjoining its fixed finite
$P_{20}$ stage stays in that family.  The compactness argument
already gives one coherent correction, rather than independent
choices at the finite stages.

The reproducible certificates are
\texttt{de\_ferrari\_a6\_kernels.py},
\texttt{de\_ferrari\_a6\_source\_data.py},
\texttt{de\_ferrari\_a6\_kernel\_source.py},
\texttt{de\_ferrari\_a6\_kernel\_corrections.py}, and
\texttt{de\_ferrari\_a6\_kernel\_compare.py}, with their
retained JSON, NPZ and stdout files.  Both modes retain the
entire correction images, and the comparison file checks the
literal quotient generators and four-strand derived groups.
The source kernels in this statement and their common
refinements pull back to normal finite-index subgroups of
$B_3$ contained in $PB_3$ after intersection with the retained
$K_6$ kernel.  No cofinality in all finite $B_3$ quotients is
asserted.
\end{proof}
